\documentclass[10pt,a4paper]{article}
\usepackage[utf8]{inputenc}
\usepackage[T1]{fontenc}
\usepackage[spanish,es-tabla]{babel}
\usepackage{csquotes}
\usepackage[margin=2cm]{geometry}
\usepackage{graphicx}
\usepackage{adjustbox}
\usepackage{booktabs}
\usepackage{tabularx}
\usepackage[backend=biber,style=ieee]{biblatex}
\usepackage[hidelinks]{hyperref}
\addbibresource{referencias.bib}

\newcommand{\cvLL}{\fontsize{17.28}{21}\selectfont}

\setlength{\parindent}{0pt}
\setlength{\parskip}{4pt}

\begin{document}
	
	% ---------- Caratula ----------
	\newgeometry{margin=2.5cm}
	\begin{titlepage}
		\begin{center}
			\fontsize{14.4}{20}\selectfont
			\vspace*{0.3cm}
			{\cvLL\bfseries UNIVERSIDAD T\'ECNICA ESTATAL DE QUEVEDO}\\[6pt]
			Facultad de Ciencias de la Computaci\'on\\[4pt]
			Carrera de Software
			\par\vfill
			
			Asignatura: Aplicaciones Distribuidas (ISR-701)\\[0.8cm]
			{\bfseries Actividad:}\\[4pt]
			FORM-AC-01 -- An\'alisis de Caso:\\
			Fallo Real en un Sistema Distribuido\\[4pt]
			Caso asignado: Amazon S3, regi\'on US-EAST-1\\
			(28 de febrero de 2017)
			\par\vfill
			
			{\bfseries Grupo GLMR -- Integrantes:}\\[6pt]
			Garc\'ia Bazurto Kevin Leonel\\
			{\fontsize{12}{16}\selectfont PFC: {[PENDIENTE]}}\\[6pt]
			Lombeida Escaleras Bryan Humberto\\
			{\fontsize{12}{16}\selectfont PFC: {[PENDIENTE]}}\\[6pt]
			Manotoa Alava Carlos Alexander\\
			{\fontsize{12}{16}\selectfont PFC: {[PENDIENTE]}}\\[6pt]
			Reinoso V\'elez Eduardo David\\
			{\fontsize{12}{16}\selectfont PFC: SSDC (Sistema de Selecci\'on Docente y Convocatorias)}
			\par\vfill
			
			Nivel: S\'eptimo\\[4pt]
			Paralelo: A\\[4pt]
			Fecha: 14 de noviembre de 2026
			\par\vfill
			
			{\bfseries Repositorio del trabajo:}\\[8pt]
			\adjustbox{max width=\linewidth}{\url{https://github.com/ereinosov/app-dist-form-ac-01}}
			\vspace*{0.3cm}
		\end{center}
	\end{titlepage}
	\restoregeometry
	
	% ---------- Contenido ----------
	\section{Descripci\'on del incidente}
	Fuente principal: informe oficial de AWS~\cite{aws-s3-2017}. {[PENDIENTE]}
	
	\begin{table}[h]
		\centering
		\small
		\begin{tabularx}{\linewidth}{@{}llX@{}}
			\toprule
			Hora (PST) & Hora (UTC) & Evento \\
			\midrule
			{[PENDIENTE]} & {[PENDIENTE]} & {[PENDIENTE]} \\
			{[PENDIENTE]} & {[PENDIENTE]} & {[PENDIENTE]} \\
			{[PENDIENTE]} & {[PENDIENTE]} & {[PENDIENTE]} \\
			{[PENDIENTE]} & {[PENDIENTE]} & {[PENDIENTE]} \\
			{[PENDIENTE]} & {[PENDIENTE]} & {[PENDIENTE]} \\
			{[PENDIENTE]} & {[PENDIENTE]} & {[PENDIENTE]} \\
			\bottomrule
		\end{tabularx}
		\caption{L\'inea de tiempo del incidente.}
	\end{table}
	
	\section{Causas ra\'iz}
	\begin{figure}[h]
		\centering
		\IfFileExists{diagramas/causas.pdf}{\includegraphics[width=\linewidth]{diagramas/causas.pdf}}{\fbox{[PENDIENTE: diagrama de causas]}}
		\caption{Diagrama de causas del incidente (elaboraci\'on propia).}
	\end{figure}
	{[PENDIENTE]}
	
	\section{Impacto cuantificado}
	{[PENDIENTE]}
	
	\section{Relaci\'on con la Unidad 1}
	\textbf{Transparencia de distribuci\'on rota.} {[PENDIENTE]}
	
	\textbf{Teorema CAP.} {[PENDIENTE]}
	
	\textbf{Tolerancia a fallos.} {[PENDIENTE]}
	
	\section{Mitigaci\'on aplicada por la empresa}
	\textbf{Medidas inmediatas.} {[PENDIENTE]}
	
	\textbf{Cambios permanentes.} {[PENDIENTE]}
	
	\section{Lecciones para el proyecto propio}
	\begin{table}[h]
		\centering
		\scriptsize
		\begin{tabularx}{\linewidth}{@{}p{2.0cm}p{2.0cm}XXXX@{}}
			\toprule
			PFC / integrante & Componente & Escenario de fallo & Relaci\'on con el caso & Medida para TA-PFC-E1 & C\'omo se verificar\'a \\
			\midrule
			{[PENDIENTE]} & {[PENDIENTE]} & {[PENDIENTE]} & {[PENDIENTE]} & {[PENDIENTE]} & {[PENDIENTE]} \\
			{[PENDIENTE]} & {[PENDIENTE]} & {[PENDIENTE]} & {[PENDIENTE]} & {[PENDIENTE]} & {[PENDIENTE]} \\
			{[PENDIENTE]} & {[PENDIENTE]} & {[PENDIENTE]} & {[PENDIENTE]} & {[PENDIENTE]} & {[PENDIENTE]} \\
			{[PENDIENTE]} & {[PENDIENTE]} & {[PENDIENTE]} & {[PENDIENTE]} & {[PENDIENTE]} & {[PENDIENTE]} \\
			\bottomrule
		\end{tabularx}
		\caption{Riesgos concretos del PFC derivados del caso.}
	\end{table}
	
	% ---------- Fuera de las 2 paginas ----------
	\clearpage
	\section*{Tabla de retroalimentaci\'on incorporada}
	\begin{tabularx}{\linewidth}{@{}p{0.6cm}p{2.7cm}XX@{}}
		\toprule
		N. & Origen & Observaci\'on recibida & Cambio realizado o raz\'on para no aplicarlo \\
		\midrule
		1 & {[PENDIENTE]} & {[PENDIENTE]} & {[PENDIENTE]} \\
		\bottomrule
	\end{tabularx}
	
	\section*{Declaraci\'on de uso de inteligencia artificial}
	\begin{tabularx}{\linewidth}{@{}p{3cm}X@{}}
		\toprule
		Herramienta y versi\'on & {[PENDIENTE]} \\
		Para qu\'e se us\'o & {[PENDIENTE]} \\
		C\'omo se verific\'o lo obtenido & {[PENDIENTE]} \\
		\bottomrule
	\end{tabularx}
	
	\printbibliography[title={Referencias}]
	
\end{document}